\chapter{Rassegna della letteratura e analisi comparativa}
\label{ch:literature-review}

\section{GNN applicati ai mercati finanziari}
\label{sec:gnn-financial-markets}

I capitoli precedenti hanno costruito gli strumenti; questo capitolo verifica cosa se ne è fatto finora e con quali risultati. È anche il capitolo in cui le due questioni di ricerca della \cref{sec:objectives-research-questions} ricevono la risposta che una tesi compilativa può dare: non una misura originale, ma una lettura critica di ciò che la letteratura ha misurato.

Conviene organizzare la rassegna per **tipo di grafo** anziché per architettura, e la ragione emerge dalla rassegna stessa: le architetture usate nei lavori applicati sono quasi sempre le stesse tre o quattro già discusse nel \Cref{ch:from-cnn-to-gnn}, mentre ciò che davvero distingue un lavoro dall'altro — e che spiega buona parte delle differenze di risultato — è la struttura relazionale scelta. La domanda progettuale non è \enquote{quale GNN}, ma \enquote{quale grafo}.

\subsection*{Grafi da relazioni societarie note}

Il filone più consolidato costruisce il grafo da relazioni tra imprese documentate indipendentemente dai prezzi. Chen et al.\ \cite{chen2018incorporating} usano relazioni tra società quotate sul mercato cinese e mostrano che una rete convoluzionale su grafo apprende rappresentazioni che catturano quelle relazioni, migliorando la previsione rispetto a modelli che trattano ogni titolo separatamente. Feng et al.\ \cite{feng2019temporal} formulano il problema come \emph{ranking} anziché come regressione — l'obiettivo è ordinare i titoli per rendimento atteso, non prevederne il valore — e costruiscono il grafo da relazioni di settore e da relazioni fornitore-cliente estratte da una base di conoscenza. Matsunaga et al.\ \cite{matsunaga2019exploring} usano relazioni di fornitura tra imprese giapponesi e combinano convoluzione su grafo e LSTM, valutando il modello con analisi a finestra mobile sul Nikkei 225.

Questi lavori condividono un vantaggio strutturale che vale la pena rendere esplicito, perché il \Cref{ch:financial-series-and-graph-construction} ne segnalava l'assenza nel caso di questa tesi: il grafo è \emph{esogeno}. Una relazione fornitore-cliente esiste indipendentemente dai prezzi, quindi il problema di circolarità della \cref{sec:dynamic-graphs-temporal-question} — la struttura stimata dagli stessi dati che si vogliono prevedere — semplicemente non si pone. In compenso, questi grafi sono statici o quasi: non catturano il cambiamento di regime che è invece il secondo oggetto di interesse di questa tesi.

\subsection*{Grafi da correlazione}

La famiglia a cui appartiene l'impianto proposto qui costruisce il grafo dalle serie dei rendimenti, con i criteri discussi nella \cref{sec:correlation-matrix-to-graph}. Il vantaggio è l'applicabilità universale — non serve alcuna fonte esterna, bastano i prezzi — e la capacità di seguire il mercato nel tempo. I costi sono quelli già esaminati: la circolarità appena citata, la fragilità statistica della stima (\cref{sec:correlation-dependence-pitfalls}) e la densità del grafo risultante (\cref{sec:known-limitations-gnn}).

Qui si trova una conferma della separazione anticipata nella \cref{sec:correlation-matrix-to-graph}: i lavori che usano il grafo di correlazione come substrato di apprendimento lo filtrano con soglie o vicinanze, mentre l'MST e il PMFG compaiono quasi esclusivamente nei lavori descrittivi di econofisica discussi nella \cref{sec:network-structure-crisis-regimes}. Le due comunità usano la stessa matrice di partenza e la filtrano in modi diversi perché rispondono a domande diverse, e raramente si citano a vicenda.

\subsection*{Grafi da testo e knowledge graph}

Una terza famiglia deriva la struttura da fonti testuali: co-menzione di due società nella stessa notizia, oppure relazioni estratte da un grafo di conoscenza. HATS \cite{kim2019hats} appartiene a questo filone e introduce un elemento rilevante per il seguito: le relazioni disponibili sono di tipo eterogeneo (stesso settore, stesso gruppo, fornitura, partecipazione azionaria) e non tutte sono ugualmente informative, per cui il modello usa un meccanismo di attenzione gerarchica per pesarle invece di trattarle alla pari. È la stessa logica che nella \cref{sec:message-passing-paradigm} rendeva il GAT interessante per il caso cripto: quando la struttura è incerta, conviene farne apprendere il peso al modello.

\begin{table}[t]
\centering
\caption[Lavori GNN su mercati finanziari per tipo di grafo]{Lavori rappresentativi di applicazione di reti neurali su grafi ai mercati finanziari, organizzati per tipo di grafo. La colonna \enquote{architettura} riporta la componente di aggregazione sul grafo; quasi tutti i lavori la combinano con un modulo temporale nel senso della \cref{sec:dynamic-graphs-temporal-question}.}
\label{tab:gnn-finance-works}
\small
\begin{tabular}{@{}llll@{}}
\toprule
\textbf{Lavoro} & \textbf{Tipo di grafo} & \textbf{Architettura} & \textbf{Compito} \\
\midrule
Chen et al.\ \cite{chen2018incorporating}  & relazioni societarie & GCN & prezzo \\
Feng et al.\ \cite{feng2019temporal}       & settore, fornitura   & GNN relazionale & ranking \\
Matsunaga et al.\ \cite{matsunaga2019exploring} & fornitura       & GCN + LSTM & prezzo \\
Kim et al.\ \cite{kim2019hats}             & relazioni eterogenee & attenzione gerarchica & direzione \\
Yu et al.\ \cite{yu2018stgcn}              & rete stradale        & GCN + conv.\ temporale & traffico \\
\bottomrule
\end{tabular}
\end{table}

La \Cref{tab:gnn-finance-works} riassume i lavori citati, includendo STGCN \cite{yu2018stgcn} come termine di paragone esterno alla finanza: la sua presenza serve a ricordare che le architetture spazio-temporali usate in questo campo nascono in domini dove il grafo è fisicamente dato, e vengono trasferite alla finanza con un cambio di ipotesi che la \cref{sec:dynamic-graphs-temporal-question} ha già segnalato.

Il quadro complessivo, per come lo restituisce anche la rassegna generale di \textcite{wu2021survey}, è che i guadagni riportati rispetto a modelli univariati esistono ma sono di entità moderata, e che la variabilità tra lavori si spiega più con il grafo e con il protocollo di valutazione — su cui la \cref{sec:evaluation-metrics-methodological-pitfalls} sarà severa — che con la sofisticazione dell'architettura. È una conclusione utile a dimensionare la risposta alla prima questione di ricerca: la struttura relazionale aiuta, ma non trasforma un problema difficile in uno facile.

\section{Applicazioni specifiche alle criptovalute}
\label{sec:applications-cryptocurrencies}

Sul mercato delle criptovalute la letteratura si divide nettamente in due: quella che ne studia la struttura di dipendenza, ampia e matura, e quella che vi applica reti neurali su grafi, ancora esigua. Vale la pena esaminarle separatamente, perché è nello scarto tra le due che si colloca l'interesse del tema.

\subsection*{Quello che si sa sulla struttura del mercato}

Il primo filone è quello richiamato nella \cref{sec:crypto-market-interconnected}. \Textcite{corbet2019cryptocurrencies} sistematizzano l'analisi delle criptovalute come classe di attività, documentandone il grado di integrazione con i mercati tradizionali. \Textcite{ji2019dynamic} studiano la connettività dinamica interna al mercato cripto con il framework di misurazione degli \emph{spillover} di \textcite{dieboldyilmaz2012spillover}, mostrando che le relazioni tra criptovalute sono intense e variabili nel tempo. I due episodi di crisi già usati come esempi nel \Cref{ch:introduction} — il collasso di Terra/Luna \cite{briola2023anatomy} e quello di FTX \cite{akyildirim2023understanding} — sono stati analizzati esattamente con questi strumenti.

Il punto rilevante è metodologico: questa letteratura descrive relazioni, spesso con notevole raffinatezza, ma raramente le usa per prevedere. Gli strumenti impiegati sono modelli econometrici di connettività o statistiche di rete calcolate a posteriori, non modelli previsivi addestrati. È il limite già osservato nella \cref{sec:why-graph-intuition}, e coincide con lo spazio che un approccio GNN potrebbe occupare.

\subsection*{Quello che si è provato a prevedere}

Le applicazioni di reti su grafi al cripto sono, allo stato, poche. Un esempio rappresentativo è il lavoro di \textcite{yin2024forecasting}, che combina convoluzione su grafo e LSTM — la stessa impostazione della \cref{sec:dynamic-graphs-temporal-question} — per prevedere il prezzo di quattro criptovalute insieme a Tether e a un indice di stress finanziario, riportando prestazioni superiori sia a una LSTM univariata sia a una LSTM multivariata su tutti gli orizzonti previsivi considerati. È un risultato coerente con la prima questione di ricerca, ma va preso per quello che è: un singolo lavoro, su pochi asset, con un protocollo di valutazione che la \cref{sec:evaluation-metrics-methodological-pitfalls} inviterebbe a esaminare con attenzione prima di generalizzare.

La scarsità di questa letteratura non va gonfiata in un merito. Significa che il campo è meno battuto di quello azionario, il che rende il tema interessante, ma significa anche che mancano i confronti sistematici — stesse criptovalute, stesso periodo, stesso protocollo — che permetterebbero di dire quanto il grafo aiuti davvero. Su questo punto la risposta onesta alla prima questione di ricerca è che l'evidenza disponibile è favorevole ma sottile.

\subsection*{Un grafo completamente diverso: le transazioni on-chain}

Esiste una terza possibilità, concettualmente distante da tutto quanto discusso finora e che merita di essere nominata perché evita alla radice il problema più fastidioso di questa tesi. Sulle blockchain pubbliche ogni transazione è osservabile, e se ne può costruire un grafo in cui i nodi sono indirizzi o transazioni e gli archi sono i flussi di valore tra essi. Il lavoro di riferimento è quello di \textcite{weber2019antimoney}, che applica reti convoluzionali su grafo a un insieme di oltre 200\,000 transazioni Bitcoin etichettate — il dataset Elliptic — per distinguere transazioni lecite da illecite.

Qui il grafo è \emph{osservato}, non stimato: nessuna finestra da scegliere, nessun coefficiente di correlazione da giustificare, nessuna circolarità tra struttura e dati da prevedere. Il prezzo di questa solidità è che il grafo risponde a domande diverse. Un grafo di indirizzi dice chi ha pagato chi, non come si muoveranno insieme i prezzi di due asset: è lo strumento adatto ad analisi forensi e antiriciclaggio, non alla previsione multi-asset. Le due rappresentazioni non sono in competizione, vivono su livelli di aggregazione differenti — l'una sui singoli partecipanti, l'altra sugli asset — e questa tesi si colloca deliberatamente sul secondo, che è il livello a cui le due questioni di ricerca sono formulate.

\section{Struttura di rete in regimi di crisi}
\label{sec:network-structure-crisis-regimes}

La seconda questione di ricerca riguarda come la topologia del grafo cambi nel passaggio da un mercato stabile a uno in crisi. È la questione su cui la letteratura offre le risposte più solide, per una ragione semplice: si tratta di misurare una struttura osservata, non di prevedere qualcosa.

\subsection*{Il fenomeno}

Il risultato ricorrente, documentato su mercati e periodi diversi, è che nelle fasi di crisi la correlazione media tra attività aumenta e la struttura del mercato si \enquote{compatta}. \Textcite{onnela2003dynamics} lo mostrano seguendo nel tempo l'albero ricoprente minimo costruito con la distanza di Mantegna della \cref{sec:correlation-matrix-to-graph}: durante i crolli l'albero si contrae topologicamente, con i titoli che si dispongono più vicini al nodo centrale invece di distribuirsi su rami lunghi. Sul versante spettrale, l'analisi delle matrici di correlazione empiriche di \textcite{laloux1999noise} — già usata nella \cref{sec:correlation-dependence-pitfalls} per separare segnale e rumore — individua un autovalore dominante, nettamente fuori dal supporto di Marchenko--Pastur dell'equazione~\eqref{eq:marchenko-pastur}, interpretabile come modo di mercato: quando quel modo assorbe una quota crescente della varianza totale, tutte le attività si muovono di fatto insieme.

Le due osservazioni sono la stessa cosa vista da due angolazioni, ed entrambe si traducono nella medesima conseguenza per il grafo: aumenta il peso medio degli archi, aumenta la densità dopo qualunque filtraggio a soglia fissa, e si riduce la separabilità del grafo in gruppi distinti.

\subsection*{Le metriche, e cosa ciascuna coglie}

Il \Cref{ch:mathematical-foundations} fornisce gli strumenti per rendere quantitative queste affermazioni. Le grandezze usate in letteratura, ordinate da quella più grossolana a quella più informativa, sono le seguenti.

\begin{description}
    \item[Correlazione media.] La media delle $\rho_{ij}$ sulle coppie. Immediata da calcolare e da interpretare, ma cieca alla struttura: due mercati con la stessa correlazione media possono avere uno un blocco fortemente coeso e il resto scollegato, l'altro una coesione uniforme.
    \item[Densità del grafo.] La frazione di archi che superano il filtraggio, nel senso della \cref{sec:weighted-graphs-notation}. Coglie il compattamento, ma dipende interamente dalla soglia scelta, e una soglia fissa rende la metrica non confrontabile tra periodi.
    \item[Lunghezza normalizzata dell'MST.] La somma delle distanze degli archi dell'albero, divisa per il numero di archi. È la metrica di \textcite{onnela2003dynamics}, robusta perché l'MST ha sempre $N-1$ archi e quindi non dipende da alcuna soglia: cala quando il mercato si compatta.
    \item[Connettività algebrica $\lambda_2$.] L'autovalore di Fiedler introdotto nella \cref{sec:spectral-analysis}. È la metrica concettualmente più adatta al fenomeno: $\lambda_2$ misura quanto sia difficile separare il grafo in due parti tagliando archi di peso ridotto, quindi cresce esattamente quando il mercato smette di articolarsi in gruppi distinti e si muove in blocco. A differenza della correlazione media tiene conto di come i pesi sono distribuiti sulla struttura, e a differenza della densità non richiede una soglia arbitraria.
    \item[Entropia spettrale.] Calcolata sulla distribuzione degli autovalori normalizzati, misura quanto la varianza sia concentrata su pochi modi: cala quando il modo di mercato domina.
\end{description}

Nessuna di queste metriche risponde da sola alla domanda, ma insieme la rendono verificabile. Questo salda la promessa fatta nella \cref{sec:spectral-analysis}, dove $\lambda_2$ era stato introdotto come strumento di analisi spettrale e il suo uso come indicatore di crisi rimandato a questo capitolo.

\subsection*{Il caso cripto}

Sul mercato delle criptovalute la letteratura è più recente ma converge. \Textcite{briola2023anatomy} analizzano il collasso di Terra/Luna documentando una riorganizzazione della rete di dipendenza tra 61 criptovalute nel periodo del crollo. \Textcite{akyildirim2023understanding} studiano il collasso di FTX con un'analisi di connettività dinamica basata su un VAR a parametri variabili nel tempo, identificando nei token emessi dalla stessa piattaforma le principali sorgenti di trasmissione degli effetti al resto del mercato. Il quadro è coerente con quello dei mercati tradizionali, con un'intensità maggiore: le criptovalute erano già fortemente correlate in condizioni normali, e nelle crisi lo diventano quasi completamente.

C'è però una tensione che questa tesi ha incontrato più volte e che qui si chiude. Il regime di crisi è il momento in cui la struttura è più interessante da studiare, ed è anche quello in cui il grafo è più denso — cioè quello in cui una GNN funziona peggio, per l'over-smoothing della \cref{sec:known-limitations-gnn}, e in cui una soglia calibrata su un periodo calmo lascia passare quasi ogni arco (\cref{sec:correlation-matrix-to-graph}). Le due questioni di ricerca non sono quindi indipendenti: il regime che rende la seconda interessante è lo stesso che rende la prima difficile. Riconoscerlo è probabilmente la conclusione più utile che questa rassegna produca.

\section{Come si valuta un modello previsivo, e come lo si valuta male}
\label{sec:evaluation-metrics-methodological-pitfalls}

Le sezioni precedenti hanno ripetutamente rimandato a questa: i risultati riportati in letteratura valgono quanto vale il protocollo con cui sono stati ottenuti, e in finanza quel protocollo è particolarmente facile da sbagliare. La sezione è breve ma non accessoria, perché stabilisce con quale grado di fiducia si possano leggere i numeri delle \cref{sec:gnn-financial-markets,sec:applications-cryptocurrencies}.

\subsection*{Quale metrica}

L'errore quadratico medio (RMSE) e l'errore assoluto medio (MAE) sono le metriche di default per la regressione, e sono quasi inutili su rendimenti finanziari. Il motivo discende direttamente dai fatti stilizzati della \cref{sec:financial-time-series-returns}: i rendimenti hanno media prossima a zero e nessuna autocorrelazione lineare, quindi il predittore costante $\hat r_t = 0$ ottiene già un RMSE vicino a quello ottimo. Un modello può migliorare l'RMSE di una frazione di punto percentuale rispetto a quel riferimento banale, e il miglioramento risulta indistinguibile dal rumore; peggio, un RMSE eccellente è perfettamente compatibile con un modello che sbaglia sistematicamente il segno.

La \emph{directional accuracy} — la frazione di volte in cui il segno previsto è corretto — misura ciò che serve per operare, ma ha un difetto speculare: tratta allo stesso modo un movimento dello $0{,}1\%$ e uno del $20\%$, per cui un modello può azzeccare il segno nel $60\%$ dei casi e perdere denaro, se sbaglia proprio nelle giornate che contano.

Le metriche economiche — rendimento cumulato di una strategia costruita sulle previsioni, Sharpe ratio, massimo drawdown — sono le uniche che rispondono alla domanda che interessa davvero, ma dipendono da ipotesi che spesso restano implicite: costi di transazione, slippage, capienza del mercato agli ordini simulati. Un backtest senza costi di transazione dichiarati non è confrontabile con nulla. La conclusione ragionevole, adottata dai lavori più solidi, è di riportare più metriche insieme e di dichiarare le ipotesi di costo.

\subsection*{Tre modi di ingannarsi}

\paragraph{Look-ahead bias.} Consiste nell'usare, per prevedere l'istante $t$, informazione non disponibile prima di $t$. La \Cref{sec:dynamic-graphs-temporal-question} aveva annunciato la forma che assume in questa tesi: la finestra su cui si stima il grafo deve chiudersi \emph{prima} dell'istante da prevedere, altrimenti la struttura relazionale contiene già il futuro. La variante più insidiosa è però un'altra, perché non ha nulla a che vedere con il grafo: standardizzare le serie usando media e deviazione standard calcolate sull'intero campione, addestramento e test insieme. È un'operazione che sembra innocua e di routine, e che inietta nel modello informazione sulla distribuzione futura dei dati.

\paragraph{Data snooping.} Se si provano cento configurazioni sullo stesso insieme di test e si riporta la migliore, il risultato riportato è quello del massimo di cento variabili casuali, non quello di un modello. Il problema è formalizzato da \textcite{white2000reality}, che propone una procedura per verificare se le prestazioni del modello migliore superino quanto atteso per puro effetto della ricerca su molte alternative. In un campo dove i modelli sono rapidi da addestrare e i dati sono pochi, il rischio è sistemico, e raramente i lavori dichiarano quante configurazioni abbiano provato.

\paragraph{Validazione errata.} La $k$-fold cross-validation, procedura standard nel machine learning, suddivide i dati in blocchi e addestra su tutti tranne uno: applicata a una serie temporale, addestra sul futuro per validare sul passato. \Textcite{bergmeir2012use} discutono in dettaglio in quali condizioni la validazione incrociata resti utilizzabile su serie temporali e in quali no. La procedura corretta nel caso previsivo è la \emph{walk-forward validation}: si addestra su una finestra iniziale, si valuta sul periodo immediatamente successivo, si sposta la finestra in avanti e si ripete — è la stessa impostazione a finestra mobile usata da \textcite{matsunaga2019exploring}, e ha il pregio di riprodurre le condizioni in cui il modello opererebbe davvero.

\subsection*{Cosa manca alla letteratura}

Chiudendo il capitolo, le lacune che emergono dalla rassegna sono tre, e sono lacune constatate, non congetturate.

La prima è quantitativa: le applicazioni di reti su grafi al mercato delle criptovalute sono poche (\cref{sec:applications-cryptocurrencies}), e mancano confronti condotti sugli stessi asset, negli stessi periodi e con lo stesso protocollo. Senza di essi, la risposta alla prima questione di ricerca resta favorevole ma sottile.

La seconda è metodologica: una parte consistente della letteratura applicata riporta miglioramenti rispetto a baseline univariate senza dichiarare costi di transazione, numero di configurazioni provate o schema di validazione. I risultati non sono necessariamente sbagliati, ma non sono verificabili, e questo suggerisce di leggere con prudenza anche i lavori favorevoli all'approccio proposto qui.

La terza è concettuale, ed è quella che questa tesi ha affrontato più direttamente. Nella maggior parte dei lavori il grafo di correlazione è assunto come dato: si sceglie una soglia, si costruisce la matrice e si passa all'architettura, senza discutere il segno dei pesi, la fragilità della stima o l'effetto della densità sul modello. I \Cref{ch:from-cnn-to-gnn,ch:financial-series-and-graph-construction} hanno mostrato che nessuna di queste scelte è neutra — l'attenzione appresa come alternativa ai pesi imposti (\cref{sec:message-passing-paradigm}), la trasformazione di Mantegna per il problema del segno e il filtraggio come condizione di funzionamento (\cref{sec:correlation-matrix-to-graph}) — ed è su questo terreno, più che sull'architettura, che una discussione più attenta avrebbe il maggiore margine di miglioramento.
